\documentclass[10pt,english]{article}

\usepackage{amsfonts,url,tikz,amsmath}

% Paper setup
\evensidemargin=0in
\oddsidemargin=0in
\textwidth=6.25in
\topmargin=-0.5in
\headheight=0.0in
\headsep=0.5in
\textheight=9.0in
\footskip=0.5in

\input{institution}

\newcommand{\Interior}[1]{\ensuremath{{#1}^{\circ}}}
\newcommand{\Closure}[1]{\ensuremath{\overline{#1}}}
\newcommand{\Complement}[1]{\ensuremath{{#1}^{c}}}

\newcommand{\Expect}{\ensuremath{\mathrm{E}}}
\newcommand{\vecnot}{\underline}
\newcommand{\RealNumbers}{\ensuremath{\mathbb{R}}}
\newcommand{\RationalNumbers}{\mathbb{Q}}
\newcommand{\ComplexNumbers}{\mathbb{C}}
\newcommand{\Real}{\mathrm{Re}}
\newcommand{\Span}{\mathrm{span}}
\newcommand{\Rank}{\mathrm{rank}}
\newcommand{\Nullity}{\mathrm{nullity}}
\newcommand{\Trace}{\mathrm{tr}}
\newcommand{\Diag}{\mathrm{diag}}
\newcommand{\dd}{\mathrm{d}}
\DeclareMathOperator*{\esssup}{ess\,sup}

% Use < , > inner product
\newcommand{\inner}[2]{{\left\langle #1 \mskip2mu , #2 \right\rangle}}
\newcommand{\tinner}[2]{{\langle #1 \mskip1mu , #2 \rangle}}

% Use < | > inner product
%\newcommand{\inner}[2]{{\left\langle #1 \mskip2mu \middle| \mskip2mu #2 \right\rangle}}
%\newcommand{\tinner}[2]{{\langle #1 \mskip1mu | \mskip1mu  #2 \rangle}}



\begin{document}

\title{ECE 586: Vector Space Methods \\ Lecture 13: Frequently Asked Questions}
\author{Henry D. Pfister \\ \univ}
\date{Revised October 19, 2026}

\maketitle

\section{Norms and Metrics}

\paragraph{What extra structure does a norm add to a vector space?}
A norm assigns a nonnegative length to each vector, is zero only at zero, scales by $|s|$, and satisfies the triangle inequality.
It induces the metric $d(v,w)=\|v-w\|$, so vector-space questions can also involve convergence and continuity.
These are the definitions on slide 1, ``Normed Vector Spaces.''

\paragraph{Why is $\min\{\|v\|_2,1\}$ not a norm even though the corresponding bounded distance is a metric?}
It fails absolute homogeneity.
For example, for a standard unit vector $e_1$, the proposed length of $2e_1$ is one, whereas twice the proposed length of $e_1$ is two.
A metric does not require compatibility with scalar multiplication; a norm does.
This is the counterexample on slide 4, ``Norms Versus Metrics.''

\paragraph{Can the exponent $p$ be less than one in a $p$-norm?}
For $p\geq1$, the usual $p$-norm satisfies the triangle inequality.
For $0<p<1$, it generally fails: in $\mathbb{R}^2$, $\|e_1+e_2\|_p=2^{1/p}>2=\|e_1\|_p+\|e_2\|_p$.
For $p=\infty$, use the largest absolute coordinate instead of a sum.

\paragraph{Why must we restrict function and sequence spaces to elements of finite norm?}
A norm must have a finite real value at every vector.
For example, $(1,1,\ldots)$ does not belong to $\ell^p$ for finite $p$, but does belong to $\ell^\infty$.
For the integral norms on slide 3, $L^p$ spaces use measurable functions with finite norm, identifying functions equal almost everywhere.
This identification is needed because an integral can ignore differences on a set of measure zero.

\paragraph{What is the difference between a supremum and an essential supremum?}
A supremum considers all point values, whereas an essential supremum ignores changes on sets of measure zero.
For continuous functions on a compact interval the two coincide with the maximum absolute value.
For general measurable functions, the essential supremum is the appropriate $L^\infty$ norm on equivalence classes.

\paragraph{Do equivalent norms assign the same length?}
No. They satisfy $m\|v\|_a\leq\|v\|_b\leq M\|v\|_a$ for fixed positive constants $m,M$.
Consequently, they have the same convergent sequences, Cauchy sequences, and topology.
All norms on a fixed finite-dimensional real or complex vector space are equivalent, as on slide 5; this need not hold in infinite dimensions.

\section{Completeness and Infinite Sums}

\paragraph{What makes a normed space a Banach space?}
Completeness in its induced metric: every Cauchy sequence must converge to a vector in the space.
For example, $C[a,b]$ with the supremum norm is Banach, because a uniform limit of continuous functions is continuous.
The same set with the $L^2$ norm need not be complete, as the example in Video 7 shows.

\paragraph{How is a Schauder basis different from a Hamel basis?}
A Hamel basis expresses every vector uniquely as a finite linear combination.
A Schauder basis allows a series whose partial sums converge in norm, with unique coefficients in the given order.
The standard unit vectors form a Schauder basis of $\ell^p$ for $1\leq p<\infty$ (slide 7), but their finite span contains only finitely supported sequences.
Not every Banach space has a Schauder basis.

\paragraph{Why does absolute convergence imply convergence in a Banach space?}
If $\sum_i\|v_i\|<\infty$, the triangle inequality gives
\[
\left\|\sum_{i=m+1}^n v_i\right\|\leq\sum_{i=m+1}^n\|v_i\|\longrightarrow0.
\]
Thus the partial sums are Cauchy, and completeness supplies their limit in the space.
This is the proof on slide 8, ``Convergence of Sums.''

\paragraph{Does being closed under linear combinations make a subspace a closed set?}
No. Algebraic closure uses finite sums, whereas topological closure also includes limits.
In $\ell^p$, the finite sums $u_n=\sum_{i=1}^n i^{-2}e_i$ converge in norm to an infinite-support sequence.
The subspace of finitely supported sequences does not contain that limit and is therefore not closed (slide 9).
Every finite-dimensional subspace of a normed space is closed.

\end{document}
